\section{Direct strict vanishing for the original quadratic-eight quartic rows}
\label{sec:cubic-tau8-strict-vanishing}

This argument concerns the original field and coefficient ring.  It uses
the already certified class sectors in AC-G135 and AC-G137 and new exact
local images of their actual global units.  No auxiliary automorphic
host, native support ideal, model condition, or coefficient ramification
restriction is used.  All strict local conditions remain literal.

\paragraph{Characters and conventions.}
Put $p=269$, $K=\mathbb Q(\sqrt{-3})$, and $t=\sqrt[3]{p}$ with its
real embedding.  Fix a generator $\chi:G_K\to\mu_3$ of $K(t)/K$.
Let $\Psi$ be the arithmetic Galois character of the quartic Dirichlet
character modulo five with $\Psi(2)=i$, choosing $i\equiv82\pmod p$.
Thus $\Psi(g)=\psi(a)$ when $g(\zeta_5)=\zeta_5^a$.
Write $c_p$ for the actual action on $p$-power roots of unity,
$\omega_p$ for its Teichmuller part, and
\[
 \tau_8=\chi_8,\qquad \tau_p=\omega_p^{134}=\chi_{269}.
 \tag{T8.1}
\]
The latter equality uses $p\equiv1\pmod4$.  Both quadratic characters
are even.  Fourth roots of unity already belong to $\mathbb Q_p$.
For any discrete coefficient module $V$, put
\[
 \operatorname{Sel}_{\rm str}(K,V)=\ker\left(
 H^1_{\rm cts}(G_K,V)\longrightarrow
 \prod_{v\mid p}H^1_{\rm cts}(G_{K_v},V)
 \times\prod_{v\nmid p,\infty}H^1_{\rm cts}(I_v,V)\right).
 \tag{T8.2}
\]
At $p$ this is zero in full local cohomology.  Away from $p$ it is
unramifiedness; no Bloch--Kato local condition replaces either one.

\begin{theorem}[AC-G141: original quadratic-eight quartic strict vanishing]
\label{AC-G141}
Let $E/\mathbb Q_p$ be finite, with its original ring of integers $O$,
uniformizer $\varpi$, and residue field $k$, containing the values of
$\chi$.  Let $\kappa:G_K\to1+\mathfrak m_O$ be any continuous character
unramified at every finite place away from $p$.  For $\sigma=\pm1$,
$r\in\mathbb Z$, $e\ge1$, and $j=1,2$, put
\[
 \psi=\Psi^\sigma\tau_8\omega_p^r\kappa,
 \qquad D=(E/O)(\chi^{-j}\psi^{-1}c_p^{-e}).
 \tag{T8.3}
\]
If $e+r\equiv267$ or $133\pmod{268}$, then
\[
                   \operatorname{Sel}_{\rm str}(K,D)=0.
 \tag{T8.4}
\]
The corresponding residual strict group is also zero for both cubic
and both quartic labels, over every finite field containing their
residual values.  Both the original $O$-length and its normalization
by $e_O=\operatorname{length}_O(O/pO)$ are zero.  No change of the
original $O$, assumption $e_O=1$, or finite-order assumption on
$\kappa$ is made.

In particular the actual original rows
$\psi=\Psi^\sigma\tau_8\tau_p^b$, $b=0,1$, are covered whenever
$e\equiv133$ or $267\pmod{268}$; the compensation changes the
effective residue by $134b$.  This strict conclusion does not use
the model/support theorem for those rows.
\end{theorem}

\begin{proof}
\proofstep{1. The exact Kummer characters and splitting fields.}
Let $q$ denote the residue of $e+r$ modulo $268$ and set
$\nu=0$ if $q=267$, $\nu=1$ if $q=133$.
For the residual character $a$ in (T8.3), the Kummer twist is
\[
 a=\bar\chi^{-j}\bar\Psi^{-\sigma}\tau_8\bar c_p^{-q},
 \qquad \eta=a/\bar c_p
       =\bar\chi^{-j}(\bar\Psi\tau_8\tau_p^\nu)^{-\sigma}.
 \tag{T8.5}
\]
Here $-(267+1)=0$ and $-(133+1)=134$ modulo $268$.
The quadratic factors are unchanged on taking an inverse, explaining
the common formula for both signs $\sigma$.

For $d=2p^\nu$ set
\[
 w=\sqrt d(\zeta_5-\zeta_5^{-1}),\quad F_d=\mathbb Q(w),
 \quad L_d=\mathbb Q(t,w),\quad M_d=KL_d.
 \tag{T8.6}
\]
In $\mathbb Q(\zeta_5,\sqrt d)$ the order-two antidiagonal
subgroup fixes $w$ and is exactly the kernel of
$\Psi\chi_d=\Psi\tau_8\tau_p^\nu$.  The quadratic field
$\mathbb Q(\sqrt d)$ is disjoint from $\mathbb Q(\zeta_5)$.
Consequently $F_d$ is the cyclic quartic field of that character and
\[
               w^4+5dw^2+5d^2=0.
 \tag{T8.7}
\]
Its only quadratic subfield is $\mathbb Q(\sqrt5)$.
The cubic splitting field $K(t)$ has only the abelian subfield $K$,
which ramifies at three whereas $F_d$ does not.  Hence
\[
 \operatorname{Gal}(M_d/\mathbb Q)=S_3\times C_4,
 \quad \Delta=\operatorname{Gal}(M_d/K)=C_3\times C_4.
 \tag{T8.8}
\]
The character $\eta$ is faithful on $\Delta$, so $M_d$ is its exact
splitting field.  In particular $\mu_p\not\subset M_d$, since
$[M_d:\mathbb Q]=24$ is not divisible by $268$.

\proofstep{2. Inverse-character descent and the ordinary unit sector.}
Let $S$ be all the $p$-places of $M=M_d$.  Write
\[
 \begin{split}
 V_{M,S}&=\{x\in M^\times:p\mid v(x)\ (v\nmid p)\}/M^{\times p},\\
 V_M&=\{x\in M^\times:p\mid v(x)\text{ at every finite }v\}/M^{\times p},\\
 \mathcal L_M&=\bigoplus_{v\mid p}M_v^\times/M_v^{\times p}.
 \end{split}\tag{T8.9}
\]
Prime-to-$p$ inflation--restriction, including the local decomposition
and inertia groups, and Kummer theory identify the residual strict
group with
\[
 e_{\eta^{-1}}\ker\left(V_{M,S}\otimes_{\mathbb F_p}k
                 \longrightarrow\mathcal L_M\otimes_{\mathbb F_p}k\right),
 \qquad \eta^{-1}=\bar\chi^j(\bar\Psi\tau_8\tau_p^\nu)^\sigma.
 \tag{T8.10}
\]
The inverse is essential: on $G_M$ the coefficient is $k(1)$, and
its natural Kummer cohomology has the additional twist $\eta$.
All finite descent quotients have order dividing twelve, so their
cohomology is exact in characteristic $p$.
Away from $p$, an unramified Kummer class has valuation divisible
by $p$.  A zero local class at a $p$-place is a local $p$th power
and has the same valuation property.  Therefore the kernel in
(T8.10) lies in $V_M\otimes k$, with all finite valuations imposed.
The exact equivariant ideal sequence is
\[
 0\longrightarrow E_M/E_M^p\longrightarrow V_M
         \longrightarrow\operatorname{Cl}(M)[p]\longrightarrow0,
 \qquad E_M=\mathcal O_M^\times.
 \tag{T8.11}
\]
It sends $(x)=\mathfrak a^p$ to $[\mathfrak a]$; its kernel and
surjectivity give the displayed sequence directly.

AC-G135 for $d=2$ and AC-G137 for $d=538$ prove that the
nontrivial cubic part of $\operatorname{Cl}(M_d)[p]$ is zero.
These are exactly these fields, including the shared-$269$ ramification
in the latter case.  Thus every class in (T8.10) comes from the
ordinary unit part of (T8.11).  No $S$-unit correction, provisional
class output, or class certificate for a different field is used.

\proofstep{3. The two simple unit modules.}
The complex places of $M$ form a free transitive $\Delta$-set after
choosing representatives over one complex embedding of $K$.
The equivariant logarithmic unit theorem gives the character of
$E_M$ modulo torsion as the regular character of $\Delta$ minus one.
Its multiplicities persist modulo $p$: the character idempotents
after an unramified scalar extension have denominators dividing
twelve and give direct summands of the integral unit lattice.
There are no $p$-power roots of unity in $M$.

Choose $b\in C_4$ with $(\Psi\tau_8\tau_p^\nu)(b)=i$ and
$\tau\in C_3$ with $\tau(t)=\zeta_3t$.  Define
\[
 e_\sigma=\tfrac14\sum_{a=0}^3i^{-\sigma a}b^a,\qquad
 e_{\rm nt}=1-\tfrac13(1+\tau+\tau^2),\qquad
 W_\sigma=e_{\rm nt}e_\sigma(E_M/E_M^p).
 \tag{T8.12}
\]
Each $W_\sigma$ has dimension two over $\mathbb F_p$.
Since $p\equiv2\pmod3$, it is the simple nontrivial
$\mathbb F_p[C_3]$-module.  Over $k$ its two cubic eigenlines
are precisely the two lines in (T8.10).
It therefore suffices to give a global unit whose localization
is nonzero on each $W_\sigma$.

\proofstep{4. The fixed global unit and its exact certificate.}
For each $d$ let $U_d$ be the fourth unit in the exact tensor-unit
certificate of AC-G135 ($d=2$) or AC-G137 ($d=538$).
These are actual algebraic elements, not unnamed numerical embeddings:
\[
 U_d=D_d^{-1}\sum_{a=0}^2\sum_{b=0}^3N^{(d)}_{ab}t^aw^b,
 \qquad D_2=3,\quad D_{538}=3p^2.
 \tag{T8.13}
\]
The twelve integers and the monic integral degree-twelve annihilator
with constant one are fixed in the accompanying exact files
\texttt{inputs/d2\_unit4\_exact.json} and
\texttt{inputs/d538\_unit4\_exact.json}.
Their SHA256 values are respectively
\begin{quote}\small\ttfamily
cbb81494267c46864b72aaf81ce559d1d\\
c8efa717446786246d5edb20314c16e\\[2pt]
aecd4d5b1fffe707587dee42096b016c\\
c3dd1fc59225f579d5b8ea1541c9a7b4.
\end{quote}
The retained class proofs already verify these unit identities.
The new finite verifier independently substitutes each annihilator
in the rational algebra with (T8.7) and $t^3=p$ and obtains zero.
Monicity and constant one prove integrality of $U_d$ and $U_d^{-1}$;
no completeness of the candidate unit list is needed.

\proofstep{5. All four unramified-quartic completions.}
For $d=2$, (T8.7) has four distinct roots modulo $p$,
\[
                           11,26,243,258.
 \tag{T8.14}
\]
Their Hensel lifts are in $\mathbb Q_p$, whereas $K_p/\mathbb Q_p$
is unramified quadratic and $t^3=p$ is Eisenstein.  Thus $M_2$
has four $p$-places, each with absolute $e=3,f=2$.
The quartic automorphism is
$b(w)=-(3+w^2/2)w$, and $b^2(w)=-w$.
Let $w_a=b^a(w)$ and put
\[
 A_a=D_2^{-1}\sum_{b=0}^3N^{(2)}_{0b}w_a^b,\qquad
 C_a=D_2^{-1}\sum_{b=0}^3N^{(2)}_{1b}w_a^b.
 \tag{T8.15}
\]
All coefficients are $p$-integral.  At each of the four roots,
$A_a\equiv1\pmod p$ for every $a$; since $p=t^3$ this gives
$b^a(U_2)=1+C_at+O(t^2)$.
At the root $w\equiv11$, the ordered $C_a$ are
$(249,134,20,135)$ modulo $p$.

Use the actual integer exponent vectors
\[
 (m_{+,a})=(202,114,67,155),\qquad
 (m_{-,a})=(202,155,67,114),
 \tag{T8.16}
\]
representing $\tfrac14i^{-\sigma a}$ modulo $p$.  The global units
\[
 A_\sigma=\prod_{a=0}^3 b^a(U_d)^{m_{\sigma,a}},\qquad
 V_\sigma=A_\sigma^3/
                 \bigl(A_\sigma\tau(A_\sigma)\tau^2(A_\sigma)\bigr)
 \tag{T8.17}
\]
satisfy $[V_\sigma]=3e_{\rm nt}e_\sigma[U_d]$.
The cubic norm has zero linear term, so their exact expansions are
\[
\begin{array}{c|rrrr}
 w\bmod p&11&26&243&258\\ \hline
 (V_+-1)/t\bmod t&166&107&162&103\\
 (V_--1)/t\bmod t&43&29&240&226
\end{array}.
 \tag{T8.18}
\]
Each entry is $3\sum_am_{\sigma,a}C_a$ modulo $p$.

\proofstep{6. Both ramified-quartic completions.}
For $d=538=2p$ put $R=w^2/p$ and $u=-w/t$.  Direct substitution gives
\[
 \begin{gathered}
 R^2+10R+20=0,\qquad u^6=pR^3,\qquad t=u^2/R,\quad w=-u^3/R,\\
 b(R)=-10-R,\qquad b(w)=(-3-R/2)w,\qquad b^2(w)=-w.
 \end{gathered}\tag{T8.19}
\]
The two simple roots of the polynomial for $R$ modulo $p$ are
$121$ and $138$.  Its Hensel lifts are $p$-adic units in
$\mathbb Q_p$; the equation for $u$ is Eisenstein of degree six.
Together with $K_p$ this gives exactly two $p$-places of $M_{538}$,
each of absolute $e=6,f=2$.

Set $B=-3-R/2$ and $(\lambda_a)=(1,B,-1,-B)$.
Replace $w$ by $\lambda_aw$ in (T8.13), substitute (T8.19),
and reduce $u^6=pR^3$.  This gives six coefficients in
$\mathbb Q[R]/(R^2+10R+20)$; the finite verifier checks that
every rational denominator in each reduced coefficient is
prime to $p$ before either residue root is evaluated.
The constant and linear coefficients are exactly
\[
 A_a=\frac{N^{(538)}_{00}+N^{(538)}_{02}\lambda_a^2pR}{3p^2},\qquad
 C_a=-\frac{N^{(538)}_{21}\lambda_ap+
                    N^{(538)}_{23}\lambda_a^3p^2R}{3p^2}.
 \tag{T8.20}
\]
Both $t^2w$ and $t^2w^3$ contributions have been retained.
The constants reduce to one, and the linear coefficients are
\[
\begin{array}{c|rr}
 a&R\equiv121&R\equiv138\\ \hline
0&109&49\\1&251&222\\2&160&220\\3&18&47
\end{array}.
 \tag{T8.21}
\]
Here $\tau(u)=\zeta_3^{-1}u$, so the same cubic norm again has
zero linear term.  The global units (T8.17) satisfy
\[
\begin{array}{c|rr}
 &R\equiv121&R\equiv138\\ \hline
 (V_+-1)/u\bmod u&91&71\\
 (V_--1)/u\bmod u&236&76
\end{array}.
 \tag{T8.22}
\]

\proofstep{7. Local injectivity and residual strict zero.}
Every displayed linear coefficient is nonzero.  If such a unit
were $y^p$, its residue one would force the residue of $y$ to be
one.  Writing $y=1+x$ in a completion of absolute ramification
degree $e_v\in\{3,6\}$ gives
\[
 v((1+x)^p-1)\ge\min(e_v+v(x),pv(x))\ge e_v+1>1.
 \tag{T8.23}
\]
This contradicts its nonzero uniformizer-linear term.  Thus
localization is nonzero on each $W_\sigma$.
The map to the sum over all $p$-places is $\Delta$-equivariant;
its kernel on the simple $\mathbb F_p[C_3]$-module $W_\sigma$
is zero.  Scalar extension preserves that injection and covers
both cubic eigenlines over $k$.
Steps 1--3 placed each residual strict class in exactly one of
these lines, with zero localization.  It must therefore vanish.

\proofstep{8. The original coefficient ring and continuous lifts.}
Write $a$ for the full character of $D$.  Its residual character
is the one already treated, since $\bar\kappa=1$.
Multiplication by the original uniformizer yields the exact sequence
\[
 0\longrightarrow k(\bar a)\longrightarrow D
         \xrightarrow{\varpi}D\longrightarrow0,
 \qquad \bar x\longmapsto x/\varpi+O.
 \tag{T8.24}
\]
At $p$ the cubic inertia has order three.  The residual rational cyclotomic
inertia has order dividing $268$, and $\Psi,\tau_8$ are unramified.
Taking the $268$th power of an inertia element with nontrivial cubic
value kills the residual rational factor and preserves the cubic
value.  Thus some $a(g)-1$ is an $O$-unit, giving
$D^{G_{K_p}}=D^{G_K}=0$.
Away from $p$, the characters $c_p,\omega_p,\kappa$ are unramified;
the remaining inertia character has finite order dividing twelve.
Its invariants are either $D$ or zero, hence are $\varpi$-divisible.
This includes the places over two, three and five.

The global long exact sequence identifies
$H^1(G_K,k(\bar a))=H^1(G_K,D)[\varpi]$.
For each full local group at $p$ and each away inertia group,
divisibility of its invariants makes the corresponding residual
$H^1$ map injective.  Hence a strict $\varpi$-torsion class lifts
to a residual strict class and is zero.
Every continuous cocycle from compact $G_K$ to discrete $E/O$
has finite image, annihilated by some power of $\varpi$.
A nonzero strict class would therefore have a last nonzero
uniformizer multiple, contradicting strict $\varpi$-torsion zero.
This proves (T8.4) over the unchanged original $O$.
\end{proof}

\paragraph{Exact scope and source dependence.}
The class statements are precisely the nontrivial cubic-sector
conclusions of AC-G135 and AC-G137.  Their finite saturation and
analytic interval certificates are retained with their established
scope; no new class-number equality or complete unit basis is claimed.
The new verifier uses exact rational arithmetic, all unit annihilators,
all four or two local embeddings, and all six ramified coefficients.
The Kummer sequence and equivariant logarithmic unit theorem are
retained as in AC-G130/AC-G134; the inverse-character descent,
all-primes ideal sequence, local injectivity and original-uniformizer
argument have been proved here.  No Lean build is asserted.

The theorem closes the original $\tau_8$ strict-zero subgate in the
specified quartic residue rows.  It neither supplies an all-prime
or growing-character average nor proves the mixed theta/Selmer
realization or the K35 common-window index bound.
Geometrically it concerns the stated Kummer torsors over
$\operatorname{Spec}\mathcal O_{M_d}[1/p]$ trivial at every
$p$-adic completion, and their original-coefficient lifts.
